\section*{集合论。}
\phantomsection
\addcontentsline{toc}{section}{集合论。}

可以说，在任何时代，数学家和哲学家都在或多或少自觉地使用着来自集合论的推理；但在他们关于这一主题的概念的历史中，必须把一切与基数观念（特别地，与无穷观念）相联系的问题，同那些仅仅引入属于与包含这两个概念的问题截然分开。后者属于最直观的概念之列，似乎从未引起过争议：正是可以在它们之上最方便地建立三段论的理论（如莱布尼茨和欧拉将向我们表明的那样），或者诸如“整体大于其各部分之和”之类的公理，更不必说几何中处理曲线与曲面相交的内容了。直到 19 世纪末，谈论由具有这样或那样给定性质的对象组成的集合（某些作者称之为“类”）都不曾出现任何问题；\(^{38}\) 而康托尔给出的那个著名的“定义”（“所谓集合，是指把我们的直觉或思维中彼此分明的一些对象汇集为一个整体”（{[}47{]}，第~282~页））在其发表之时几乎没有引起任何异议。\(^{39}\) 然而一旦集合的概念与数或量的概念掺和在一起，情况就完全不同了。直线的无限可分性问题（最早的一批毕达哥拉斯学派学者无疑已经提出了它）众所周知后来引出了相当大的哲学困难：从埃利亚学派到波尔查诺和康托尔，数学家和哲学家们前仆后继地扑向“由无穷多个无大小的点组成的有限的量”这一悖论，却均未成功。即使概略地重述这一问题所激起的无休止而激烈的论战——它们构成一块特别有利于形而上学或神学离题发挥的土壤——对我们来说也没有任何益处；我们只需指出自古以来大多数数学家所止步的那个观点。这一观点实质上在于：只要还不能以无可辩驳的方式作出结论，就拒绝讨论——我们将在现代形式主义者那里重新见到这种态度；正如后者设法排除“悖论的”集合的介入一样（见下文第~31~页以下），经典数学家也小心翼翼地避免在推理中引入“实无穷”（也就是说，避免引入由被设想为同时存在的、至少在思想中同时存在的无穷多个对象所组成的集合），而满足于“潜无穷”，即增大每个给定量的可能性（若是“连续”的量，则是减小它的可能性）。\(^{40}\) 尽管这一观点含有一定分量的虚伪，\(^{41}\) 它却毕竟容许了经典数学的绝大部分（包括比例理论以及后来的无穷小演算）得到发展；\(^{42}\) 尤其是在无穷小所引起的论战之后，它甚至被视为一种绝妙的保险，直到深入 19 世纪之后很久，它几乎仍是人们普遍信奉的一条教条。

关于等势这一一般概念的第一缕微光，出现在伽利略的一段评注中（{[}122 b{]}，第 VIII 卷，第~78-80~页）：他注意到映射 \(n \mapsto n^2\) 在自然数与其平方数之间建立了一个双射对应，并由此得出结论：公理“整体大于部分”

不能应用于无穷集合。但是，这一评注远没有开创对无穷集合的理性研究，它似乎除了加深人们对实无穷的不信任之外别无其他效果；这已经是伽利略本人所得出的结论，而柯西在 1833 年引用他，也只是为了赞同他的这一态度。

正是分析的需要——特别是贯穿整个 19 世纪的对实变量函数的深入研究——构成了后来成为现代集合论的那门理论的源头。当波尔查诺在 1817 年证明 R 中一个有下界的集合的下极限存在时（{[}27 c{]}），他仍像他同时代的大多数人那样“按内涵”推理，谈论的不是任意的实数集合，而是实数的任意性质。但 30 年后，当他撰写《无穷的悖论》（Paradoxien des Unendlichen）{[}27 b{]}（1851 年出版，即他去世三年之后）时，他已毫不犹豫地主张“实无穷”存在的权利，并谈论任意的集合了。他在这部著作中定义了两个集合等势的一般概念，并证明了 R 中两个紧区间是等势的；他还注意到有限集与无穷集的特征区别在于：无穷集 E 与一个异于 E 的子集等势；但他对这一断言没有给出任何令人信服的证明。这部著作的总基调终究是哲学多于数学，而且由于未能把集合的势的概念同无穷的大小或等级的概念足够清楚地分离开来，波尔查诺构造势越来越大的无穷集的尝试均告失败，他还借此机会把推理与关于发散级数的、完全没有任何意义的考虑混在了一起。

我们今天所知的集合论之得以创立，应归功于 G. 康托尔的天才。他同样是从分析出发的：他关于三角级数的工作（受黎曼工作的启发）很自然地使他在 1872 年对出现在这一理论中的“例外”集合作出第一次分类尝试，\(^{43}\) 所借助的则是他在此时引入的逐次“导集”的概念。无疑正是与这些研究、也与他定义实数的方法（{[}47{]}，第~92-97~页）相联系，康托尔开始对等势问题发生兴趣，因为在 1873 年他指出，有理数集（或代数数集）是可数的；而在大约始于此时与戴德金的通信 {[}48{]} 中，我们看到他提出了整数集与实数集之间的等势问题，几周之后他便成功地给出了否定的解决。随后，从 1874 年起，占据他心灵的是维数问题：有三年的时间，他徒劳地试图证明 R 与 \(R^{n}(n > 1)\) 之间不可能存在双射对应，最终却大出其所料地\(^{44}\)定义出了这样一个对应。有了这些既新奇又惊人的结果，他便把自己完全献身于集合论。在 1878 年至 1884 年间发表于《数学年刊》（Mathematische Annalen）的一系列 6 篇论文中，他同时处理了等势问题、全序集理论、R 与 R \(^{n}\) 的拓扑性质以及测度问题；令人赞叹的是，可以看到这些在经典“连续统”观念看来如此错综纠缠的概念，在他手中是如何一点一点地清晰呈现出来的。从 1880 年起，他有了“超限地”迭代“导集”构造的想法；但这一想法直到两年之后，随着良序集——康托尔最独创的发现之一——的引入才获得生命；多亏良序集，他才能够着手对基数作详细的研究，并提出“连续统问题”{[}47{]}。

这样大胆的观念推翻了两千年的传统，又导出如此出人意料、外表如此悖理的结果，它们不可能不经抵抗就被接受。事实上，在当时德国有影响的数学家中，只有魏尔斯特拉斯一人以几分好感关注着康托尔（他昔日学生）的工作；另一方面，康托尔却遇到了施瓦茨尤其是克罗内克的不可调和的反对。\(^{45}\) 看来，使康托尔出现神经疾病的最初症状、并使他的数学创造力因之受损的，既是由对其思想的反对所造成的持续紧张，同样也是他为证明连续统假设而付出的徒劳努力。\(^{46}\) 他直到 1887 年才真正重新拾起对集合论的兴趣，他最后的论著发表于 1895-1897 年；其中他主要发展了全序集理论和序数的演算。他还在 1890 年证明了不等式 \(m < 2^{m}\)；然而，不仅连续统问题仍无答案，而且在基数论中还留下一个更严重的空缺：康托尔未能确立任意基数之间良序的存在。这一空缺即将得到填补：一方面靠 F. 伯恩斯坦定理（1897），它表明关系 \(a \leq b\) 与 \(b \leq a\) 蕴涵 \(a = b\)，\(^{47}\) 另一方面尤其是靠策梅洛定理 {[}342 a{]}，它证明了每个集合上都存在良序——这个定理康托尔早在 1883 年就已猜测到了（{[}47{]}，第~169~页）。

然而戴德金从一开始便以持续的关注追随着康托尔的研究；不过，当后者把注意力集中于无穷集合及其分类时，戴德金却循着他自己关于数概念的思路（这条思路此前已使他借“分割”给出了无理数的定义）。在他 1888 年出版的小册子《数是什么，应是什么》（Was sind und was sollen die Zahlen）中——其要旨早在 1872-78 年即已写成（{[}79{]}，第 III 卷，第~335~页）——他表明了自然数概念（如前所见，整个经典数学最终都依赖于它）本身如何能够从集合论的基本概念中导出。他（无疑是最早明确地这样做的人）发展了一个集合到另一个集合的任意映射的初等性质（这在康托尔之前一直无人问津，康托尔只对双射对应感兴趣）；对于集合 E 到自身的每个映射 f，他引入了元素 \(a \in E\) 相对于 f 的“链”的概念，即集合 \(K \subset E\) 中所有满足 \(a \in K\) 且 \(f(K) \subset K\) 者的交。⁴⁸ 随后，他把无穷集 E 定义为：存在 E 到 E 内的一个双射 \(\phi\) 使得 \(\phi(E) \neq E\)；⁴⁹ 进一步，若还存在这样一个映射 \(\phi\) 和一个元素 \(a \notin \phi(E)\)，使得 E 是 a 的链，戴德金便称 E 是“单纯无穷的”；他注意到此时“皮亚诺公理”即得到满足，并（先于皮亚诺）表明了如何从出发得到初等算术的全部定理。他的阐述中唯一缺失的部分是无穷公理：戴德金（步波尔查诺的后尘）自认为通过把人的“思想世界”（“Gedankenwelt”）看作一个集合就能证明这条公理。⁵⁰

从另一条路走，戴德金由于其在算术方面的工作（特别是理想理论）而被引去在比康托尔更一般的框架中考虑有序集的概念。当后者把自己完全局限于全序集时，\(^{51}\) 戴德金则处理一般的情形，并对格作了特别深入的研究（{[}79{]}，第 II 卷，第~236-271~页）。这项工作在当时几乎无人注意；尽管这些被好几位作者重新发现的结果自 1935 年以来成为大量出版物的对象，它们的历史重要性与其说来自这一理论无疑相当有限的应用可能，不如说来自如下事实：它们构成了精心进行公理化构造的最早例子之一。另一方面，康托尔关于可数集或具有连续统的势的集合的最初结果，很快就有了为数众多而且重要的应用，直至延伸到分析中最经典的问题\(^{52}\)（自然还不必说康托尔工作中开创了一般拓扑学和测度论的那些部分；关于这一点见第~141~页和第~221~页）。此外，从 19 世纪的最后几年起，超限归纳原理的最初应用开始出现；尤其在策梅洛定理获证之后，它已成为现代数学所有分支中不可或缺的工具。库拉托夫斯基于 1922 年给出了这条原理的一个往往更合用的版本，它避免了良序集的使用（{[}189 a{]}，第~89~页）；今天人们主要使用的正是这个形式，它后来又由佐恩重新发现 {[}344{]}。\(^{53}\)

到 19 世纪末，康托尔的基本概念就这样站住了脚。\(^{54}\) 我们已经看到，大约在这个时候，数学的形式化宣告完成，公理化方法的使用也几乎取得了一致的同意。但与此同时，一场罕见的剧烈的“基础危机”开启了，它震撼数学界达三十余年，而且有时看来不仅危及一切新近的收获，甚至危及数学中最经典的部分。

\section*{集合论的悖论与基础危机。}
\phantomsection
\addcontentsline{toc}{section}{集合论的悖论与基础危机。}

最初的“悖论的”集合出现在基数和序数的理论中。1897 年，布拉利-福尔蒂指出，不能考虑由全部序数组成的集合的存在，因为这个集合将是良序的，从而同构于它的一个异于全体的段，而这是荒谬的 {[}42{]}。\(^{55}\) 1899 年，康托尔（在一封致戴德金的信中）注意到，无论谈论基数组成一个集合，还是谈论“所有集合的集合”，都不再可能不陷入矛盾（这最后一个“集合” \(\Omega\) 的子集的集合将与 \(\Omega\) 的一个子集等势，这与不等式 \(m < 2^{m}\) 相悖）（{[}47{]}，第~444-448~页）。最后，1905 年，罗素在分析这个不等式的证明时表明，建立该不等式的推理（无须求助于基数论）同时也证明了“所有不是其自身元素的集合所组成的集合”这一概念是矛盾的 {[}266{]}。\(^{56}\)

人们或许会以为，这样的“二律背反”只会出现在数学的边缘地带，即出现在以考虑直觉无法企及的“大小”的集合为特征的区域。然而另一些“悖论”很快就威胁到数学中最经典的部分。贝里和罗素（{[}266{]}，第 I 卷，第~63-64~页）简化了 J. 里夏尔的一个论证 {[}258{]}，他们实际上注意到：定义可用英语少于十六个词表述的整数所组成的集合是有限的，然而把一个整数定义为“用英语少于十六个词所不能定义的最小整数”却是矛盾的，因为这个定义本身只由十五个词组成。

这类推理与数学家通常使用的推理相距甚远，在许多人看来也许只像一种字谜，但它们毕竟指出了修改数学基础、以期消除这类“悖论”的必要性。然而，尽管关于这种修改的紧迫性人人意见一致，关于实现它的方式，很快就出现了根本的分歧。对第一派数学家——无论他们是“观念论者”还是“形式主义者”\(^{57}\)——来说，集合论“悖论”所造成的局面，与几何学中因发现非欧几何或“病态”曲线（如无切线的曲线）而产生的局面十分类似；由此应当得出一个类似但更为一般的结论，即想把任意一门数学理论建立在（明确的或不明确的）对“直觉”的诉诸之上，是毫无希望的。这一立场可以用形式主义学派主要论敌的原话来概括：“形式主义者”，布劳威尔说（{[}39 a{]}，第~83~页），“认为人的理性并不拥有例如直线或大于十的数的精确图像\ldots 诚然，对于数学实体之间的某些关系，即我们当作公理的那些关系，我们可以按照固定的规则从其他关系把它们推导出来，并相信这样一来我们便是通过逻辑推理从其他真理导出了真理\ldots{[}但是{]}对形式主义者来说，数学的精确性只存在于关系串的发展过程之中，它独立于人们可能愿意赋予这些关系或它们所联结的实体的任何意义。”

因此，对形式主义者来说，事情就在于给集合论一个与初等几何十分类似的公理基础：在那里，人们不必关心被称为“集合”的那些“对象”是什么，也不必关心关系 \(x \in y\) 意味着什么，而只需列举出应加于这个关系上的条件；当然，这样做时必须尽可能把康托尔理论的全部结果都包括进来，同时又使“悖论的”集合不可能存在。这种公理化的第一个例子由策梅洛于 1908 年给出 {[}342 c{]}；他通过引入“分离公理（Aussonderung）”来避开“过大的”集合：性质 \(P(x)\) 仅当 \(P(x)\) 已蕴涵某个形如 \(x \in A\) 的关系时，才确定出一个由具有该性质的元素组成的集合。⁵⁸ 但是，要消除类似于“里夏尔悖论”的悖论，只能靠限制附加在“性质”概念上的意义来实现；在这一点上，策梅洛只满足于用十分含糊的措辞描述一类他称之为“已定义的”性质，并指出在应用分离公理时必须限于这些性质。这一点由斯科朗 {[}286 a{]} 和弗兰克尔 {[}113 b{]} 加以精确化；正如他们所指出的，对这一点的澄清迫使我们置身于一个完全形式化的系统之中，在那里，“性质”与“关系”这些概念已失去全部“意义”，只是按照明确的规则构造出来的汇集的简单称谓。这当然就要求把所使用的逻辑规则也放进系统之内，而这在策梅洛-弗兰克尔系统中尚付阙如。

在此之后，人们又提出了集合论的其他一些公理化。其中主要应提到冯·诺伊曼的公理化 {[}342 a and c{]}，它比策梅洛-弗兰克尔系统更接近康托尔的原初构想：后者为避开悖论的集合，早在与戴德金的通信中就已提议（{[}47{]}，第~445-448~页）区分两类集合，即“杂多”（“Vielheiten”）与本来意义上的“集合”（“Mengen”），后者的特征在于它们可以被设想为单一的对象。冯·诺伊曼正是把这一想法精确化，区分出两类对象：“集合”与“类”；在他的系统中（几乎是完全形式化的），类与集合的区别在于：类不能被置于记号 ∈ 的左边。这种系统的好处之一，是它为 19 世纪逻辑学家所使用的“全类”概念恢复了名誉（全类自然不是集合）；还应指出，冯·诺伊曼的系统（就集合论而言）避免了公理模式的引入，而代之以合适的公理（这使得对它作逻辑研究更为容易）。贝尔奈斯和哥德尔给出了冯·诺伊曼系统的一些变体 {[}130 b{]}。

上述系统看来已经妥善地消除了悖论，但代价是引入了一些只能让人觉得十分任意的限制。为策梅洛-弗兰克尔系统说句公道话：它所作的限制只不过认可了把集合概念应用于各门数学理论时的现行实践。冯·诺伊曼和哥德尔的系统离通常的观念更远；但另一方面，某些尚处于初创阶段的数学理论，要安插在这样的系统所提供的框架中，说不定反而比安插在策梅洛-弗兰克尔系统较窄的框架中更容易，这并非不可能。

当然，谁也不能断言这些解决方案中的任何一个给人以定论的印象。它们之所以让形式主义者满意，是因为后者拒绝考虑每个数学家的个人心理反应；他们假定，当一种形式化语言能够把数学推理转录为一种没有歧义、从而适于充当数学思想之载体的形式时，它便完成了自己的任务；他们会说，只要一个人的推理能够转录进共同的语言，那么他对数学对象的“本性”或他所使用的定理的“真理性”爱怎么想就怎么想。\(^{59}\)

换句话说，从哲学的观点看，形式主义者的态度就在于对“悖论”所提出的问题不闻不问，从而抛弃了那种力图赋予数学概念以全体数学家共同的理智“内容”的柏拉图式立场。面对这一与传统的决裂，许多数学家望而却步。例如罗素就试图通过对悖论的结构作更深入的分析来避免它们。罗素和怀特海重新拾起由 J. 里夏尔最先提出（在他展示其“悖论”的那篇文章 {[}258{]} 中）、后又由 H. 庞加莱加以发展 {[}251 c{]} 的一个想法，他们注意到，悖论集合的定义全都违反如下这条称为“恶性循环原则”的原则：“其定义涉及一个集合的元素之全体的那种元素，不能属于该集合”（{[}266{]}，第 I 卷，第~40~页）。也正是这一陈述充当了《数学原理》（Principia）的基础，该书正是为了满足它才发展出“类型论”。与启发它的弗雷格的逻辑一样，罗素和怀特海的逻辑含有“命题变元”；类型论对这些不同的变元进行分类，其主要轮廓如下。从一个未加确定的“个体域”出发——其中的个体可称为“0 阶对象”——，凡变元（自由的或约束的）皆为个体的那些关系，称为“一阶对象”；一般地，凡变元为 ≤ n 阶对象（其中至少一个为 n 阶）的那些关系，称为“\(n + 1\) 阶对象”。\(^{60}\) 于是，n 阶对象组成的一个集合便只能由一个 \(n + 1\) 阶的关系来定义，这一条件使悖论集合得以毫无困难地被消除。\(^{61}\) 但“类型层级”的原则限制性极强，若严格坚持它，得到的就是一种复杂得无法梳理的数学。\(^{62}\) 为了摆脱这一后果，罗素和怀特海不得不引入一条“可化归性公理”，它断言：对于“个体”之间的每一个关系，都存在一个与之等价的一阶关系；这条公理与形式主义者的公理同样任意，而且大大削弱了《数学原理》这一构造的意义。因此，罗素和怀特海的系统在逻辑学家那里比在数学家那里更成功；无论如何，它并未完全形式化，\(^{63}\) 因而存在许多晦暗不明之处。人们曾作出各种努力来简化和澄清这一系统（拉姆齐、赫维斯泰克、蒯因、罗塞尔）；这些作者倾向于使用越来越完全形式化的语言，他们用一些只考虑所论聚合体的措辞的限制，取代了《数学原理》的那些规则（后者尚有某种直观的基础）；这些规则不仅看起来与策梅洛-弗兰克尔系统或冯·诺伊曼系统中所表述的禁令同样毫无根据，而且由于离数学家的实践更远，它们在许多情形下导致了其创立者未曾预见的不可接受的后果（例如布拉利-福尔蒂悖论，或对选择公理的否定）。

对前述各派的数学家来说，首要之点在于不放弃过去遗产的任何部分：“谁也不能把我们从康托尔为我们创造的乐园中驱逐出去”，希尔伯特说（{[}163 c{]}，第~274~页）。为此，他们愿意接受对数学推理的限制；这些限制由于与习惯相合而无伤大体，但它们看来并非由我们的心智习惯和集合概念的直观观念所强加。在他们看来，任何事情都好过让心理学混入数学有效性的判据；与其像阿达马所说的那样“把我们大脑的特性考虑在内”（{[}12{]}，第~270~页），他们宁可听任给数学领域强加一些大体上任意的界限，只要这些界限把经典数学包括在内，又不至于妨碍未来的进步。

而当我们要看的是我们行将谈及的那一倾向的数学家的态度时，情况就完全不同了。如果形式主义者甘愿放弃对数学推理的“心灵之眼”的监督，那么被称为“经验主义者”、“实在论者”或“直觉主义者”的数学家则拒绝这一弃权；他们必须有一种内在的确定性，为保证他们所处理的数学对象的“存在”作保。当问题只在于放弃空间直观时，还没有什么严重的异议，因为算术“模型”容许人们退守到整数这个直观概念后面。但当要把整数概念化归为（直观上远不那么清晰的）集合概念、进而给集合的操作加上毫无直观基础的屏障时，不可调和的对立就显露出来了。这些反对者中最早的一位（而且凭其天才的权威，也是影响最大的一位）是 H. 庞加莱；他不仅承认几何学的公理学观点和分析的算术化，还承认康托尔理论的相当一部分（他是最早在其工作中富有成效地应用这一理论的人之一），但他却拒绝设想算术本身也能通过公理化的处理得到论证；在他看来，完全归纳原理尤其是我们心智的一种基本直觉，不可能把它看成一种纯粹的约定\(^{64}\) {[}251 c{]}。他从原则上敌视形式化语言，并质疑其用处；他不断地把形式化数学中的整数概念，与当时刚刚兴起、我们将于后文谈到的证明论中对整数的使用混为一谈；毫无疑问，在那个时代，要像今天这样——经过 50 年的研究和讨论之后——把这一区别弄清楚是很困难的，尽管一个希尔伯特或一个罗素当时已清楚地感到了它。

这类批评在策梅洛于 1904 年引入选择公理 {[}342 a{]} 之后便成倍增加。在此以前，它在分析或集合论的许多证明中的使用几乎无人觉察；\(^{65}\) 正是遵循埃哈德·施密特提示的一个想法，策梅洛把这条公理明确地表述出来，并以灵巧的方式由它导出了每个集合上都存在良序的一个令人满意的证明。由于与“悖论”同时出现，这种新的推理方式凭其古怪的姿态看来使许多数学家陷入惶惑；只要看看随后一卷《数学年刊》上署名发表的种种奇怪的误解即可，署名者竟是像舍恩弗利斯和 F. 伯恩斯坦那样熟悉康托尔方法的人。E. 博雷尔发表在同一卷上的批评更有分量，而且清楚地附和了 H. 庞加莱关于整数所表达的观点；这些批评在 E. 博雷尔、贝尔、阿达马和勒贝格——他们仍置身于法国数学的经典传统之中——的通信往来中得到展开和讨论 {[}12{]}。博雷尔一开始就否认选择公理的有效性，因为它一般包含不可数无穷多个选择，这在直观上是无法设想的。但阿达马和勒贝格指出，可数无穷多个任意的相继选择也并不更直观，因为它由无穷多次运算组成，而这些运算不可能被设想为实际发生。对把讨论范围扩大了的勒贝格来说，一切都归结为弄清说一个数学对象“存在”是什么意思；在他看来，必须“指名”一个以唯一方式定义该对象的性质（我们会说这是一种“泛函式的”性质）；对于策梅洛在其推理中所使用的那类函数，勒贝格称之为选择的“规律”；他继续说，如果这一要求得不到满足，如果人们只是“想到”这个函数而不“指名”它，那么能肯定在推理过程中人们想到的一直是同一个东西吗（{[}12{]}，第~267~页）？无论如何，这给勒贝格带来了新的疑虑；在集合中选取单个元素的问题在他看来就已引起困难：必须确信这样的元素“存在”，也就是说，能“指名”该集合的至少一个元素。\(^{66}\) 那么，对于一个不知道如何“指名”其每个元素的集合，还能谈论它的“存在”吗？贝尔已经毫不犹豫地否认一个给定无穷集合的所有子集组成的集合的“存在”（出处同上，第~263-264~页）；阿达马指出这些要求将导致放弃谈论实数的可能性，也是徒劳：E. 博雷尔最终赞同的正是这一结论。且不说可数性看来已经取得了永佃权，人们几乎是回到了“实无穷”的反对者的经典立场上。

所有这些异议都谈不上系统；要按类似但更为激进的原则对数学作一次彻底的改造，这任务落到了布劳威尔及其学派的身上。像直觉主义这样一门一半是心理学一半是数学的复杂学说，我们不敢在此加以概述，我们只限于指出其若干最引人注目的方面，欲知详情，请参看布劳威尔本人的著作 {[}39 a and b{]} 以及海廷的阐述 {[}162{]}。对布劳威尔来说，数学与我们思想中的“精确”部分同一，它奠基于自然数序列这一原初直观之上，而且不可能在不遭损毁的情况下把它翻译成一个形式系统。无论如何，它只在数学家的思想中才是“精确”的，想在他们之间发展出一种不受语言的一切缺陷和歧义之累的交流工具，纯属幻想；至多只能希望借助或多或少含糊的描述，在对话者心中唤起一种有利的心态（{[}162{]}，第~11-13~页）。直觉主义数学对逻辑的重视几乎不超过对语言的重视：一个证明之令人信服，并非凭借一劳永逸地固定的逻辑规则，而是凭借其每一步的“直接明证性”。这种“明证性”还必须以比 E. 博雷尔及其追随者更受限制的方式来解释：因此在直觉主义数学中，不能说形如“R 或（非 R）”的关系为真（排中原理），除非对 R 中变元的每一组给定的值，都能证明 R、“非 R”两个命题之一为真；例如，从两个实数之间的方程 ab = 0 不能推出“a = 0 或 b = 0”，因为很容易举出实数 a、b 的明确例子，使得 ab = 0 成立，而在当前却无法证明 a = 0、b = 0 两个命题中的任何一个（{[}162{]}，第~21~页）。

从这样的原则出发，直觉主义数学家最终得到与经典定理大不相同的结果，这不足为怪。这些结果中有整整一部分消失了，例如分析中大部分“存在性的”定理（如实函数的波尔查诺定理和魏尔斯特拉斯定理）；如果一个实变量函数在直觉主义的意义上“存在”，它便理所当然地连续；有界的实数单调序列不一定有极限。另一方面，许多经典概念在直觉主义者那里分岔为若干个不同的基本概念：例如（实数序列的）收敛概念有两个，可数性概念有八个。不言而喻，超限归纳法及其对现代分析的应用（如同康托尔理论的大部分内容一样）被不加申诉地判了死刑。

依照布劳威尔的看法，数学命题只有以这种方式才能获得“内容”；超出直觉主义者容许范围的形式主义推理被判为毫无价值，因为再也无法赋予它们一种可以应用“真”这一直观概念的“意义”了。显然，这样的判断只能依靠一种先已存在的、具有心理学或形而上学本性的“真理”概念。这实际上等于说，它们是拒绝一切讨论的。

毫无疑问，来自直觉主义阵营的猛烈攻击有时不仅迫使数学中的前卫学派、甚至迫使传统数学的追随者采取守势。有一位著名数学家承认，这些攻击给他留下如此深刻的印象，以致他自愿把工作限制在被判定为“可靠”的数学分支之内。但这类情况很难是经常发生的。直觉主义学派的下场，其记忆无疑注定只作为一种历史奇观而留存；但它至少立过一功：它迫使它的论敌——也就是数学家的压倒多数——明确了自己的立场，并更清楚地认识到他们信赖数学的理由（一部分属于逻辑方面，另一部分属于情感方面）。

\section*{元数学}
\phantomsection
\addcontentsline{toc}{section}{元数学}

无矛盾性在任何时代都被视为全部数学的一个必要条件，而且从亚里士多德的时代起，逻辑就已充分发展，人们完全清楚地知道：从一个矛盾的理论可以推出任何东西。自古代以来就被视为不可或缺的那些“存在性”证明，显然别无其他目的，只是要保证新概念的引入不至于带来矛盾，尤其是当这个概念过于复杂、无法直接纳入“直观”之下时。我们已经看到，随着 19 世纪公理学观点的到来，这一要求变得如何迫切，而算术“模型”的构造又是如何回应它的。但是，算术本身难道不会是矛盾的吗？在 19 世纪末之前，人们无疑不会想到这样的问题，因为整数显得属于我们直观中最可靠的那一部分；但在“悖论”之后，一切都似乎成了问题，于是可以理解，它们所造成的不安全感促使数学家们从 1900 年前后起更加慎重地对待算术的无矛盾性问题，以便至少使经典数学免遭灭顶之灾。这个问题因此成为希尔伯特在其著名的 1900 年国际数学家大会演讲中所列举的第二个问题（{[}163 a{]}，第 III 卷，第~229-301~页）。在这样做时，他提出了一条将产生巨大反响的新原则：在传统逻辑中，一个概念的无矛盾性只是使它成为“可能的”，而对希尔伯特来说（至少对公理地定义的数学概念而言），无矛盾性与概念的存在是一回事。这显然意味着，必须先证明一门数学理论的无矛盾性，然后才能合法地发展它；H. 庞加莱正是这样理解的，当他在形式主义身上打开缺口时，他把希尔伯特的这一想法接了过来，并怀着狡黠的快意强调，形式主义者在当时离能够做到这一点还相去甚远（{[}251 c{]}，第~163~页）。我们将在后文看到希尔伯特如何应战；但在那之前，这里必须指出，在这两位的权威的影响下，庞加莱所提出的这些要求在很长时期内被形式主义者及其论敌双方毫无保留地接受了。其后果之一是形式主义者中流传甚广的一个信念：希尔伯特的证明论是数学的一个组成部分，而且是其不可或缺的绪论。这一教条在我们看来并无根据，\(^{67}\) 我们认为，元数学对逻辑和数学的阐述的介入，可以而且应当缩减到涉及缩写记号的运用和演绎判据的那一点非常初等的部分。因此，与庞加莱的想法相反，问题不在于“要求免于矛盾”，而在于同阿达马一起认为：无矛盾性即使未被证明，也是被注意到的（{[}12{]}，第~270~页）。

剩下要做的，是对希尔伯特及其学派的努力作一简短的历史概述，并且概略地追述两条线索：一是最终导出哥德尔否定性结果、并使阿达马的怀疑态度得到事后印证的那段演变，二是此后围绕数学推理机制的认识所取得的一切进步——正是这些进步使现代元数学成为一门自有其不容置疑的意义的独立科学。

早在 1904 年，希尔伯特就在国际数学家大会的一次演讲（{[}163 c{]}，第~247-261~页）中着手处理算术的一致性问题。他首先指出，利用模型来证明它是不可能的，\(^{68}\) 并概述了另一种方法的原理：他提议把形式化算术的真命题看作没有意义的符号汇集，并证明：在使用支配这些汇集的构成与联结的规则时，不可能得到一个既是真命题、其否定又是真命题的汇集。他甚至对一个比算术更窄的形式勾勒了这种性质的证明；但是，正如 H. 庞加莱不久之后所指出的（{[}251 c{]}，第~185~页），这个证明本质上使用了递归原理，因而看来依赖于一个恶性循环。希尔伯特没有立即回应这一批评，此后大约十五年无人尝试发展他的思想；直到 1917 年（出于回应直觉主义者攻击的愿望），他才重新套上数学基础问题这驾马车，而且从此直到科学生涯结束不曾停止考虑这个问题。在他大约从 1920 年延续到 1930 年的关于这一主题的工作中——一整派年轻数学家（阿克曼、贝尔奈斯、埃尔布朗、冯·诺伊曼）参与了这些工作——希尔伯特一点一点地更精确地提出了他的“证明论”的诸原则：他默认了庞加莱批评的正确性，承认在元数学中，所使用的算术推理只能建立在我们对整数的直观（而不是形式化算术）之上；为此，看来必须把这种推理限制在直觉主义者所容许的那类“有限程序”（“finite Prozesse”）之内：例如，一个反证法证明不能证明一个汇集或汇集序列在元数学上的存在，必须有一个明确的构造规律。\(^{69}\) 另一方面，希尔伯特在两个方向上扩大了他最初的纲领：他不仅处理算术的一致性，还力图证明实数论甚至集合论的一致性；\(^{70}\) 此外，在一致性问题之外又添上了公理的独立性、完备性和判定问题。我们将逐一简述这些不同的问题，并指出它们所引出的主要研究。

命题组 \(A_{1}, A_{2}, \ldots, A_{n}\) 的独立性证明在于表明：对每个指标 i，\(A_{i}\) 在理论 \(T_{i}\) 中不是定理，该理论是以 \(A_{j}\)（其指标 \(j \neq i\)）为公理而得到的。如果知道一个一致的理论 \(T_{i}^{\prime}\)，在其中 \(A_{j}(j \neq i)\) 都是定理，而且“非 \(A_{i}\)”也是定理，这一点便得到了确立；于是，根据是否承认某些理论（如算术或集合论）是一致的，可以以两种方式考虑这个问题。在第二种情形，所面对的是一个“绝对”一致性问题。相反，第一类问题通过构造适当的“模型”化归为“相对”一致性问题，而且早在数学取得完全形式化的面貌之前，人们就已想出了许多这种类型的证明：只需回想非欧几何的模型（见第~133~页）、希尔伯特在《几何基础》（Grundlagen der Geometrie）中处理的初等几何公理的独立性问题 {[}163 c{]}，以及施泰尼茨关于代数公理化的工作和豪斯多夫及其后继者关于拓扑学公理化的工作，即足以说明。

称一个理论 T 为完备的，如果对 T 的每个命题 A（不含 T 的常元以外的任何字母），命题 A 与（非 A）之一是 \(T.^{71}\) 的定理。若把某些非常初等的形式体系撇在一旁——它们的完备性是容易证明的（{[}164{]}，第~35~页）——这个领域所得的结果本质上是否定性的；其中最早的一个属于 K. 哥德尔 {[}130 a{]}，他表明：如果 T 是一致的，且形式化算术的公理都是 T 的定理，那么 T 不是完备的。他那个巧妙方法的根本想法，在于（当然借助于“有限程序”）在元数学陈述与形式化算术中的某些命题之间建立一个双射对应；我们只限于勾画一个粗略的轮廓。\(^{72}\) 对 T 的每个作为项或关系的汇集 A，先以一个可以近乎机械地应用的（属于明确构造的）程序，以双射的方式关联一个整数 \(g(A)\)。同样，对 T 的每个证明 D（被看作汇集的一个相继序列），可以以双射的方式关联一个整数 \(h(D)\)。最后，可以给出一个明确地构造 T 的关系 \(P(x,y,z)\) 的程序，\(^{73}\) 使得在 T 中，\(P(x,y,z)\) 要求 x,y,z 是整数，并且满足以下两个条件：

\(1^{\circ}\) 若 \(D\) 是 \(A(\lambda)\) 的一个证明，其中 \(A(x)\) 是 \(\mathcal{T}\) 的一个关系，而 \(\lambda\) 是一个明确的整数（也就是说，是 \(\mathcal{T}\) 的一个是整数的项），则 \(P(\lambda, g(A(x)), h(D))\) 是 \(\mathcal{T}\) 的定理；

\(2^{\circ}\) 若明确的整数 \(\mu\) 不具有 \(h(D)\) 的形式，或者 \(\mu = h(D)\) 而 \(D\) 不是 \(A(\lambda)\) 的证明，则（非 \(P(\lambda, g(A(x)), \mu)\)）是 \(\mathcal{T}\) 的定理。

现在设 \(S(x)\) 是关系（非 \((\exists z)P(x,y,z)\)），并设 \(\gamma = g(S(x))\)，它是 \(\mathcal{T}\) 的一个项。如果 \(\mathcal{T}\) 是一致的，那么命题 \(S(\gamma)\) 在 \(\mathcal{T}\) 中就没有证明。事实上，若 \(D\) 是这样一个证明，则 \(P(\gamma,g(S(x)),h(D))\) 将是 \(\mathcal{T}\) 的定理；但这个关系不是别的，正是 \(P(\gamma,\gamma,h(D))\)，由此推知 \((\exists z)(P(\gamma,\gamma,z)\) 也将是 \(\mathcal{T}\) 的定理；由于这最后一个关系等价于（非 \(S(\gamma)\)），\(\mathcal{T}\) 就会是不一致的。另一方面，刚才所说的表明，对每个明确的整数 \(\mu\)，（非 \(P(\gamma,\gamma,\mu)\)）都是 \(\mathcal{T}\) 的定理。其结果是，\(\mathcal{T}\) 中不存在（非 \(S(\gamma)\)）的证明，因为这最后一个关系等价于 \((\exists z)P(\gamma,\gamma,z)\)，而一旦存在整数 \(\mu\) 使 \(P(\gamma,\gamma,\mu)\) 成立，依据前文所述就将要求 \(\mathcal{T}\) 是不一致的。\(^{74}\) 哥德尔的这条元数学定理随后在若干方向上得到了推广（{[}181{]}，第 XI 章）。\(^{75}\)

T 的关系 \(S(\gamma)\)——我们刚才表明，无论这个关系还是它的否定，在 T 中都没有证明——显然是专为达成目的而拼造出来的，它同任何数学问题都没有自然的联系。有意思得多的是如下事实：若以 \(\mathcal{T}\) 表示集合论（采用冯·诺伊曼-贝尔奈斯公理系统），则连续统假设及其否定在 \(\mathcal{T}\) 中都不可证明。这一著名结果是分两步建立的：1940 年，哥德尔证明了通过向 \(\mathcal{T}\) 添加连续统假设 \(2^{\aleph_0} = \aleph_1\) 而得到的理论并非不一致 {[}130 b{]}；随后在 1963 年，P. 科恩证明了：当向 \(\mathcal{T}\) 添加关系 \(2^{\aleph_0} = \aleph_2\)（或者，对任意整数取 \(2^{\aleph_0} = \aleph_n\)，其中 \(n > 1\)）时，结论同样成立 {[}67{]}。

判定问题（“Entscheidungsproblem”）无疑是元数学中所陈述的一切问题里雄心最大的一个：它在于弄清，对一个给定的形式语言，能否设想一种近乎机械的“普适程序”，当把它应用于所论形式体系的任何关系时，能在有限次运算之内指出该关系是否为真。\(^{76}\) 这一问题的解决本来就属于莱布尼茨的宏伟设想的一部分，而且希尔伯特学派有一度似乎认为其实现已近在眼前。事实上，对只含少数原始符号和公理的形式体系，人们能够描述出这样的程序（{[}181{]}，第~136-141~页）。但是，为使判定问题精确化、即确切勾勒“普适程序”意指什么而作的种种努力，迄今为止只得到了否定的结果（{[}181{]}，第~432-439~页）。此外，一个理论 T 的判定问题的解决，立即使人得知 T 是否一致，因为只需把“普适程序”应用于 T 的一个关系及其否定即可；而我们将看到，对通常的数学理论，不可能以这种方式解决这一问题。\(^{77}\)

的确，正是在数学理论的一致性问题——元数学的源头和心脏——上，结果显得最令人失望。在 1920-1930 年间，希尔伯特及其学派为攻克这些问题发展了新的方法；在证明了覆盖算术之一部分的部分形式体系的一致性之后（参见 {[}158{]}、{[}324 d{]}），他们以为自己已经达到目标，证明了不仅是算术的一致性、而且还有集合论的一致性；这时哥德尔依据算术的不完备性推出：用希尔伯特的“有限程序”不可能证明任何包含算术的理论 T 的一致性。\(^{78}\)

然而，哥德尔定理并没有把证明一致性的尝试之门完全关死，只要（至少部分地）放弃希尔伯特关于“有限程序”的限制即可。于是根岑在 1936 年 {[}127{]} 成功地证明了形式化算术的一致性，他所使用的是“直观地”超限归纳到可数序数 \(\varepsilon_{0}\) 的做法。⁷⁹ 可以加于这种推理之上的“确定性”的价值，无疑不如满足希尔伯特最初要求的推理那样令人信服，而且它本质上是每个数学家的个人心理问题；但同样真实的是：这类“直观地”超限归纳到某一给定序数的类似“证明”，若能应用于例如实数理论或集合论的一个实质部分，将会被视为重要的进步。

\chapter{记号；组合分析。}

历史学和考古学使我们了解到为数众多的“记数体系”；它们最初的目标，是给每个单独的整数（直至一个视应用需要而定的界限）配一个名称和一种书面表示，后者由为数有限的符号按或多或少规则的法则组合而成。最为常见的做法是把整数分解为“相继单位”\(b_{1}, b_{2}, \ldots, b_{n}, \ldots\) 之和，其中每一个都是前一个的整数倍；虽然一般取 \(b_{n}/b_{n-1}\) 等于同一个数 b（体系的“基”，最常见的是 10），但这一规则有不少例外，如巴比伦人那里 \(b_{n}/b_{n-1}\) 有时等于 10，有时等于 6 {[}232{]}，又如玛雅人的历法体系中，\(b_{n}/b_{n-1}\) 除 n = 2 外都等于 20，而 \(b_{2}/b_{1} = 18\) {[}228{]}。至于相应的书面表示，它必须指明每一阶 i 的“单位”\(b_{i}\) 的数目；在许多体系中（如埃及人、希腊人和罗马人那里），相继的倍数 \(k.b_{i}\)（其中 k 从 1 变到 \((b_{i+1}/b_{i}) - 1\)）由一些同时依赖于 k 和 i 的符号来表示。第一个重要的改进在于：用同一个符号表示所有的数 \(k.b_{i}\)（对 k 的同一个值）；这就是“位值记数法”的原理，其中指标 i 由如下事实来指示：表示 \(k.b_{i}\) 的符号出现在构成所表示之数的诸“切片”序列中的第 i 个位置上。这种类型的第一个体系见于巴比伦人那里：无疑早在公元前 2000 年，他们就用同一个记号表示与指数 i 的一切值相应的所有倍数 \(k.60^{\pm i}\)（{[}232{]}，第~93-109~页）。这种体系的不便之处当然在于所用符号的歧义性：只要没有任何东西指明某一阶的单位可能空缺，换句话说，只要该体系未通过引入“零”而加以补全，歧义就一直存在。然而人们看到，巴比伦人在其历史的大部分时间里并不使用这样的符号；事实上，他们只是在公元前最后两个世纪才使用这样一种“零”，而且还只在数的内部使用；在那之前，只能靠上下文来确定所论之数的意义。只有另外两个体系系统地使用“零”：玛雅人的体系（据推测在公元纪年开始时即已在使用 {[}228{]}），以及我们现今的十进制体系，它（经由阿拉伯人）源自印度数学，在印度，其使用早在公元最初几个世纪即有明证。还应指出，把零当作一个数（而不只是一个简单的分隔记号）的概念及其引入运算，也属于印度人的独创性贡献（{[}78{]}，第 I 卷）。当然，一旦掌握了位值记数法原理，把它推广到任何基都是容易的事；关于 17 世纪以来提出的各种“基”孰优孰劣的讨论属于数值计算的技术问题，不宜在此展开。我们只限于指出，作为这些体系之基础的那个运算，即所谓“欧几里得”除法，在希腊人之前并未出现，它无疑应追溯到最早的一批毕达哥拉斯学派学者，他们把它作为其理论算术的基本工具（见第~85~页）。

冠以“组合分析”之名的那些一般的计数问题，看来直到古典古代的最后几个世纪才得到处理：只有公式 \(\binom{n}{2}=n(n-1)/2\) 在公元 3 世纪得到确证。印度数学家婆什迦罗（12 世纪）知道 \(\binom{n}{p}\) 的一般公式。更系统的研究见于 13 世纪初利维·本·热尔松的一部手稿：他得到了使 n 个对象中每次取 p 个的排列数 \(V_{n}^{p}\) 得以计算的递推公式，特别是 n 个对象的排列数；他还陈述了等价于关系 \(\binom{n}{p}=V_{n}^{p}/p!\) 与 \(\binom{n}{n-p}=\binom{n}{p}\) 的法则（{[}311{]}，第 VI 卷，第~64-65~页）。但这部手稿似乎被同时代人所忽视，其结果只是在后来的几个世纪中才被数学家们逐渐重新发现。在后起的进展中，我们可以指出：卡尔丹证明了 n 个元素的集合的非空子集数是 \(2^{n}-1\)；帕斯卡和费马在创立概率演算时重新发现了 \(\binom{n}{p}\) 的表达式，而帕斯卡是第一个注意到这些数与二项式公式之间关系的人：后者看来早在 13 世纪已为阿拉伯人所知，14 世纪为中国人所知，并于 16 世纪初在西方被重新发现，同时被重新发现的还有通常归于帕斯卡名下的、称为“算术三角形”的递推计算方法（{[}311{]}，第 VI 卷，第~35-38~页）。最后，莱布尼茨在 1676 年前后得到了（但未发表）“多项式系数”的一般公式，该公式由德·莫弗独立地重新发现，并于 20 年后发表。

\chapter{代数的演进。}

在数学中，很少有比合成法则更原始的概念：它似乎与整数和可计量的量的计算的最初萌芽不可分割。我们从埃及人和巴比伦人数学留存下来的最古老的文献可以看出，他们已经掌握了一整套关于整数 \textgreater{} 0、有理数 \textgreater{} 0、长度和面积的计算法则；尽管流传到我们手中的巴比伦文本只处理其中数据取明确数值的问题，\(^{1}\) 它们却毫不掩饰地表明，所使用的法则被赋予了一般性，并且在处理一次和二次方程时显示出相当可观的技术娴熟（{[}232{]}，第~179~页及以下）。无论如何，在其中找不到对所用法则的任何论证的迹象，甚至找不到对所出现运算的精确定义：这一切连同那些法则都停留在经验主义的领域之内。

另一方面，类似的关切在古典时代的希腊人那里已经表现得十分清楚；诚然，那时还见不到整数理论的一种公理化处理（这样的公理化直到 19 世纪末才会出现；见第~24~页）；但在欧几里得的《原本》中，已有许多段落对一些其直观程度与整数计算法则相当的“明显”计算法则（例如两个有理数乘积的交换性）给出形式的证明。这类证明中最值得注意的，是与量的理论相关的那些——量的理论是希腊数学最独创的创造（如我们所知，它等价于我们的实数 \textgreater{} 0 的理论；见第~148~页以下）：欧几里得在其中考虑了两个量的比之积等对象，证明了它不依赖于给出这两个比的形式（这是合成法则对等价关系的“商”的第一个例子），并证明了它是交换的（{[}107{]}，第 V 卷，命题 22-23）。\(^{2}\)

不过我们不必讳言，在欧几里得那里，这一朝向严格性的进步，就代数计算的技术而言，是伴随着停滞、甚至某些倒退的。几何学压倒性的优势（量的理论显然正是在几何学的框架内构思出来的）瘫痪了代数记号的一切自主发展：计算中出现的元素必须时时刻刻被“几何地”表示出来；更何况，所出现的两个合成法则不是定义在同一个集合上的（比的加法没有一般性的定义，而两长度之积不是长度而是面积）；由此造成的不灵活，使得处理二次以上的代数关系几乎无法实行。

只是随着古典希腊数学的衰落，人们才看到丢番图回到“计算师”即职业计算者的传统，这些人一直在照原样应用从埃及人和巴比伦人那里继承来的法则：由于不再为他所考虑的那些“数”寻求几何表示，他很自然地被引向发展抽象代数计算的法则；例如，他给出了一些（用现代语言说）等价于公式 \(x^{m+n} = x^{m}x^{n}\)（对 m、n 的小的（正或负）值）的法则（{[}91 a{]}，第 I 卷，第~8-13~页）；稍后一点（第~12-13~页）可以找到“符号法则”的陈述，这是负数计算的最初迹象；最后，丢番图第一次使用一个文字符号来表示方程中的未知数。另一方面，他似乎几乎不去想把他用以解决问题的方法与一般观念联系起来；至于合成法则的公理化概念——它在欧几里得那里已初露端倪——无论对丢番图的思想还是对其直接后继者的思想来说都是陌生的；在代数中，要到 19 世纪初才会重新找到它。

首先需要的是，在其间的几个世纪中，一方面发展出一套适合表达抽象法则的代数记号体系，另一方面，“数”的概念得到拓宽，使人们能够通过对足够多样化的特殊情形的观察上升到一般概念。在这一点上，希腊人创造的量的比的公理化理论是不够的，因为它只是把实数 \textgreater{} 0 的直观概念以及对这些数的运算精确化了，而巴比伦人早已以更含混的形式知道了它们；现在要谈的相反是希腊人未曾设想过的那些“数”，而且起初没有任何可感的“表示”自行呈现出来：一方面是零和负数，它们早在中世纪前期就已出现于印度数学；另一方面是虚数，它们是 16 世纪意大利代数学家的创造。

如果把零撇开——它最初是作为记号中的一个符号被引入的，后来才被当作一个数（见第~45~页）——这些扩张的共同特征是（至少起初）纯粹“形式的”。这应当理解为：新的“数”最初是作为施行于某些情形的运算结果而出现的，在这些情形下，按其严格的定义，它们没有任何意义（例如当 a \textless{} b 时两个整数之差 a - b）：因此人们赋予了它们“假的”、“虚构的”、“荒谬的”、“不可能的”、“想像的”数等等名称。对于一味追求清晰思维的古典时代的希腊人来说，这样的扩张是不可设想的；它们只能出自这样一些计算者：他们比希腊人更倾向于对其方法的力量怀有甚至是几分神秘的信心（用 18 世纪的说法，“分析的普遍性”），任凭自己被计算的机制拖着走而不逐步检验其可靠性；而这种信心事后大多得到了辩护，因为把原本仅对先前已知的数严格有效的计算法则推广到这些新数学对象上去，导出了正确的结果。这就解释了，人们何以越发大胆地就其本身（独立于对具体计算的一切应用）来考虑数概念的这些推广——起初它们只是运算序列中的中介，其出发点和目标都是真正的“数”；一旦迈出这一步，人们便开始为这些新实体寻找或多或少具体的解释，它们由此获得了在数学中被引用的权利。\(^{4}\)

在这一点上，印度人已经意识到负数在某些情形下应有的解释（例如商业问题中的一笔债务）。在其后的几个世纪中，随着希腊和印度数学的方法与成果（经由阿拉伯人）传向西方，人们对这些数的使用越来越熟练，并开始有了对它们的另外一些几何的或动力学的“表示”。无论如何，这与代数记号的逐步改进一起，构成了中世纪末期代数中唯一值得注意的进步。

16 世纪初，代数开始了新的腾飞，这要归功于意大利学派的数学家发现了三次方程、随后四次方程的“根式解”（见第~72~页以下）；正是在这个场合，他们尽管心怀厌恶，却可以说是被迫把虚数引入了计算；不过，对这类“不可能的”数的计算的信心逐渐建立起来，如同对负数一样，尽管在两个多世纪里没有人想出对它们的一种“表示”。

另一方面，代数记号从韦达和笛卡尔那里获得了决定性的改进；从后者开始，代数记号除去少数例外，已经就是我们今天所使用的记号了。

从 17 世纪中叶到 18 世纪末，无穷小演算的创立所敞开的广阔视野，看来在某种程度上造成了对代数总的忽视，特别是对合成法则、对实数和复数本性的数学思考的忽视。\(^{5}\) 于是，力学中早在 17 世纪末就已熟知的力的合成与速度的合成，对代数没有产生任何影响，尽管它们已经以胚胎的形式包含着向量演算。的确，我们要一直等到 1800 年前后导致复数的几何表示的那场思想运动（见第~161~页），才在纯粹数学中看到向量加法的使用。\(^{6}\)

正是在这个时候，合成法则的概念第一次在代数中沿两个不同的方向得到扩展，扩展到这样一些元素：它们与“数”（这个词至此被赋予的最广意义）之间只剩遥远的类比。第一次扩展应归功于 C. F. 高斯，出自他关于整系数二次型 \(ax^{2} + bxy + cy^{2}\) 的数论研究。拉格朗日曾在同一判别式的型的集合中定义了一个等价关系，\(^{7}\) 并且另外证明了一个恒等式，它在该集合中给出一个交换的（并非处处有定义的）合成法则；从这些结果出发，高斯证明这个法则与一个稍有不同的等价关系相容（{[}124 a{]}，第 I 卷，第~272~页）：“由此可以看出”，他当时说，“由两个或多个类合成的类应当理解为什么”。随后他进而研究他刚刚定义的“商”法则，实质上确立了它（用现代语言说）是一条阿贝尔群法则，而且他所用的论证在多数情形下远比高斯所考虑的特殊情形更具普遍性（例如，他用来证明对称元唯一性的论证可以应用于任何合成法则（同上，第~273~页））。但他没有止步于此：稍后重回这一问题时，他认识到类的合成与整数模一个素数的乘法之间的类似\(^{8}\)（同上，第~371~页），但也注意到给定判别式的二次型的类群并不总是循环群；他在这方面给出的指示使我们想到，他或许至少在这一特殊情形下已经认识到了有限阿贝尔群的一般结构（{[}124 a{]}，第 I 卷，第~374~页和第 II 卷，第~266~页）。

我们要谈的另一系列研究同样以群的概念告终，从而以决定性的方式把它引入数学：这就是“置换理论”，它是拉格朗日、范德蒙德和高斯关于代数方程求解的思想的发展。我们不打算在这里给出这一问题的详细历史（见第~75~页以下）；我们应当记住鲁菲尼 {[}265{]}、随后是柯西（{[}56 a{]}，(2)，第 I 卷，第~64~页）对有限集的两个置换之“积”的定义，\(^{9}\) 以及关于有限置换群的最初概念：传递性、本原性、中性元、交换的元素，等等。但这些最初的研究整体上仍相当表面，而真正应当被视为这一理论开创者的是埃瓦里斯特·伽罗瓦：在他那部令人难忘的著作 {[}123{]} 中，他把代数方程的研究化归为他与这些方程相联系的置换群的研究，并对后者作了重大的深化，既涉及群的一般性质（正是伽罗瓦第一个定义了正规子群的概念并认识到它的重要性），也涉及具有特殊性质的群的确定（他在这一方面所得的结果至今仍跻身于该理论最精深的结果之列）。同样也是

伽罗瓦最先有了“群的线性表示”这一想法，\(^{10}\) 而这一事实清楚地表明，他已掌握两个群结构同构的概念，而与其“实现”无关。

然而，即使高斯和伽罗瓦的杰出方法无疑已把他们引向了十分宽广的合成法则概念，他们却未曾有机会专门展开他们在这一点上的想法，其工作对抽象代数的演进也没有产生任何直接的作用。\(^{11}\) 朝向抽象的最明显的进展是在第三个方向上取得的：继关于虚数本性的思考（其几何表示在 19 世纪初曾引起相当可观的工作）之后，英国学派的代数学家从 1830 年到 1850 年首先提炼出合成法则的抽象概念，并立即把这概念应用于一大堆新的数学对象，从而扩大了代数的领域：布尔那里的逻辑代数（见第~8~页）、哈密顿那里的向量、四元数和一般的超复系统 {[}145 a{]}、凯莱那里的矩阵和非结合法则（{[}58{]}，第 I 卷，第~127~页和第~301~页，及第 II 卷，第~185~页和第~475~页）。一场平行的演变在欧洲大陆独立发生，尤其在向量计算（默比乌斯、贝拉维蒂斯）、线性代数和超复系统（格拉斯曼）方面（见第~61~页）。\(^{12}\)

从使代数在 19 世纪上半叶重新焕发活力的这一切独创而多产的思想的沸腾之中，代数脱颖而出，连其趋向都为之一新。在此之前，方法和结果都围绕着代数方程（或数论中的丢番图方程）求解这一中心问题运转：“代数”，塞尔在其《高等代数教程》的引言中说 {[}284{]}，“严格说来就是方程的分析”。1850 年以后，尽管代数教材在很长时间里仍把首要地位留给方程理论，新的研究已不再为求数值方程之解的立竿见影的应用所支配，而是越来越转向我们今天视为代数本质问题的东西，即对代数结构本身的研究。

这项工作相当清晰地分成三股潮流，它们分别延伸着我们上面分析过的三场思想运动，并且彼此平行地推进、相互间没有可察觉的影响，直至 19 世纪的最后几年。\(^{13}\)

首先是德国学派（狄利克雷、库默尔、克罗内克、戴德金、希尔伯特）对代数数理论的建构，它脱胎于高斯的工作——最早的一类研究，即对数 \(a + bi\)（\(a\)、\(b\) 为有理数）的研究，正应归功于高斯。我们不打算在此追述这一理论的演变：我们只需为了自己的目的拣出在其中浮现的抽象代数概念。早在高斯的第一批后继者那里，（代数数的）域的概念就已处于上述全部工作的基础地位（阿贝尔和伽罗瓦关于代数方程的研究也是如此）；当戴德金和韦伯 {[}80{]} 把单变量代数函数论移植到代数数论上时，其应用范围进一步扩大。理想概念的引入也应归功于戴德金（{[}79{]}，第 III 卷，第~251~页），它给出了元素集合之间合成法则的一个新例子；阿贝尔群和模在代数域理论中所起的越来越大的作用，应追溯到他和克罗内克；我们将在后文回到这一点（见第~63~页以下和第~93~页以下）。

我们把线性代数（第~57~页）和超复系统（第~117~页）发展史的介绍也留到以后；它们在 19 世纪末和 20 世纪初的展开没有引入任何新的代数概念：在英国（西尔维斯特、W. 克利福德）和美国（B. 与 C. S. 皮尔斯、迪克森、韦德伯恩），是沿着哈密顿和凯莱开辟的道路；在德国（魏尔斯特拉斯、戴德金、弗罗贝尼乌斯、莫利恩）和法国（拉盖尔、E. 嘉当），则以独立于盎格鲁-撒克逊人的方式并采用颇为不同的方法。

至于群论，它起初主要以有限置换群理论的形式得到发展，这是在伽罗瓦著作发表、并经塞尔的著作 {[}284{]} 传播之后，尤其要归功于 C. 若尔当的巨著《置换论》{[}174 a{]}。若尔当在此书中对前人关于置换群特殊性质（传递性、本原性等）的工作作了极大的完善和总结，所得结果中的大部分至今几乎未被超越；他还深入研究了非常重要的特殊群，即线性群及其子群（见第~138~页）；同样，正是他引入了一个群由另一个群来表示这一基本概念，以及（稍后一点）商群的概念，并且证明了以“若尔当-赫尔德定理”之名著称的定理的一部分。\(^{14}\) 最后，无穷群的最早研究也应归功于若尔当 {[}174 b{]}，几年之后，S. 李与 F. 克莱因和 H. 庞加莱分别沿两个不同的方向对它作了重大的发展。

在此期间，人们逐渐看清：群的本质在于其合成法则，而不在于组成群的那些对象的性质（例如见 {[}58{]}，第 II 卷，第~123~页和第~131~页，以及 {[}79{]}，第 III 卷，第~439~页）。然而，甚至关于有限抽象群的研究，在很长时间里仍被当作关于置换群的研究来构思，直到 1880 年前后，才有了有限群自主理论的自觉发展。我们无法进一步追述这一理论的历史；我们只限于提及至今仍在有限群研究中位居最常用工具之列、而且都可上溯到 19 世纪的两种工具：其一是西罗 \(^{15}\) 关于 p 群的诸定理，其年代为 1872 年 {[}303{]}；其二是由弗罗贝尼乌斯在本世纪最后几年创立的特征标理论（{[}119{]}，第 III 卷，第~1-37~页）。愿意深入了解有限群理论及其所引出的众多困难问题的读者，可参阅专著 {[}44 b{]}、{[}131{]}、{[}291{]} 和 {[}341{]}。

同样是在 1880 年前后，群的表现开始得到系统的研究；在此之前，人们研究的只是特殊群的表现，例如交错群 \(\mathfrak{A}_5\) 在哈密顿著作 {[}145 b{]} 中的表现，或微分方程与黎曼曲面的单值群的表现（施瓦茨、克莱因、施莱夫利）。W. 迪克第一个 {[}97{]} 定义了（但未命名）由有限多个生成元生成的自由群，但他对它本身的兴趣，不如把它当作一个“普适”对象的兴趣——该对象使“由生成元和关系给出的群”这一说法获得确切的定义。迪克发展了凯莱最先提出的一个想法，而且显然还受到上文提到的黎曼学派工作的影响，他描述了具有给定表现的群的一种解释，其中每个生成元由关于两个相切圆或相交圆的反演之积来表示（另见 {[}44 b{]}，第 XVIII 章）。稍后，在庞加莱及其后继者发展了自守函数理论、庞加莱又引入了代数拓扑的工具之后，基本群的首批研究将与群表现的研究并肩推进（德恩、蒂策），两种理论相互支撑。而且恰恰是一位拓扑学家 J. 尼尔森，于 1924 年在其性质的第一个深入研究 {[}234{]} 中引入了“自由群”这一名称；紧接着，E. 阿廷（仍然与拓扑问题相关）引入了群的自由积概念，O. 施赖尔则更一般地定义了带融合子群的自由积。这里我们同样无法处理这一方向后续发展的历史，欲知详情，请参看著作 {[}216{]}。

这里也不是谈论群的概念（以及与之紧密相连的不变量概念）自 19 世纪末以来在分析、几何、力学和理论物理中所享有的异常繁荣的地方。正是以类似的方式，这一概念以及与之相关的一些代数概念（带算子群、环、理想、模）侵入了代数中直到那时看来还相当远离其本域的那些部分——我们在此追述的演化的最后一个时期即以此著称，而且它以我们上文考察过的三种倾向的合流而告终。这一统一首先是现代德国学派的功绩：代数的公理化工作发端于 19 世纪末的戴德金和希尔伯特，E. 施泰尼茨 {[}294 a{]} 强有力地推进了它，随后从 1920 年起，又在 E. 阿廷、E. 诺特及其学派的代数学家（哈塞、克鲁尔、O. 施赖尔、范德瓦尔登）的推动下继续进行。范德瓦尔登出版于 1930 年的教程 {[}317 a{]} 第一次把这些成果汇集在一个完整的阐述之中，为抽象代数此后多种多样的研究开辟了道路并充当了向导。
